\documentclass[11pt]{article}
\usepackage[margin=1in]{geometry}
\usepackage{times}
\usepackage{hyperref}
\usepackage{listings}
\usepackage{xcolor}
\usepackage{booktabs}
\usepackage{amsmath}
\usepackage{amssymb}
\usepackage{enumitem}
\usepackage{pifont}

\hypersetup{colorlinks=true, linkcolor=blue, urlcolor=blue, citecolor=blue}

\lstset{
  basicstyle=\ttfamily\small,
  breaklines=true,
  frame=single,
  columns=fullflexible,
  backgroundcolor=\color{gray!8}
}

\title{Toward a Taxonomy of Vulnerability Patterns in \\ AI-Assistant-Generated Code}
\author{Jitendra Rout \\ \small Independent Researcher \\ \small \texttt{[email/GitHub link]}}
\date{September 2026}

\begin{document}
\maketitle

\begin{abstract}
AI coding assistants have become a routine part of software development, and
a growing body of anecdotal and industry commentary suggests they introduce
security vulnerabilities in patterns distinct from those found in
human-authored code. This paper proposes a preliminary taxonomy of seven such
patterns, grounded in a proposed causal account of why each pattern arises
from the incentives of code-sample generation rather than from model
capability limits. We implement a proof-of-concept static analysis scanner,
\texttt{ai-vuln-scan}, targeting these patterns, and report an empirical
comparison against two established open-source static analysis tools
(Bandit for Python, and the \texttt{eslint-plugin-security} plugin for
JavaScript) on a small hand-constructed test corpus. We find that general-purpose
security linters, run with their default offline rule sets, miss a
substantial share of the patterns in our corpus --- most notably a class of
inter-route authorization-consistency bug with no direct analogue in either
baseline tool. We report these results candidly, including a case where our
own tool underperformed the established baseline, and outline the taxonomy,
tooling, and evaluation methodology as an open, extensible artifact for
further community contribution.
\end{abstract}

\section{Introduction}

Since the widespread adoption of AI coding assistants (GitHub Copilot,
Claude Code, and comparable tools), a recurring concern in developer and
security communities is that AI-generated code exhibits security defects at
a rate, and in a distribution, that differs meaningfully from human-authored
code. This concern is largely anecdotal in the public discourse --- security
newsletters and conference talks routinely cite examples --- but there is
comparatively little open, structured tooling built specifically around the
premise that these are \emph{distinct, nameable} patterns rather than a
generic increase in defect rate.

This paper makes three contributions:

\begin{enumerate}
  \item A taxonomy of seven vulnerability patterns (Section~\ref{sec:taxonomy}), each
  with a proposed causal account of \emph{why} the pattern is
  disproportionately associated with AI-assisted generation, not merely
  possible in any code.
  \item An open-source reference implementation, \texttt{ai-vuln-scan}
  (Section~\ref{sec:tool}), that detects instances of six of the seven
  patterns using structural and pattern-based heuristics.
  \item An empirical comparison (Section~\ref{sec:evaluation}) against two
  established, independently-maintained static analysis tools on a small
  test corpus, reported with full honesty about cases where our tool
  under- or over-performed the baseline.
\end{enumerate}

We emphasize at the outset that this is preliminary, small-sample work --
a proposed taxonomy and a proof-of-concept scanner, not a large-scale
empirical study of real-world AI-generated code. We see the primary
contribution as the taxonomy and the open, versioned, contribution-friendly
structure in which it is published, rather than a claim of comprehensive
tool performance.

\section{Related Work}

General-purpose static analysis tools for security --- Bandit
\cite{bandit} for Python, \texttt{eslint-plugin-security}
\cite{eslintsecurity} for JavaScript, and commercial/hosted tools such as
Semgrep \cite{semgrep} --- detect a broad range of vulnerability classes
using pattern matching, taint analysis, and (for some tools) hosted rule
registries. These tools are not designed around any particular hypothesis
about \emph{generation source}; a hardcoded secret is flagged identically
whether written by a human or an AI assistant. Our work is
complementary rather than competitive: we do not claim to detect a broader
class of vulnerabilities than these tools in general, only that a subset of
patterns is worth naming and tuning for specifically because of their
disproportionate association with AI-assisted generation.

Industry discussion of "AI-introduced" vulnerabilities is active but
largely takes the form of blog posts, vendor marketing content, and
conference talks rather than a shared, citable taxonomy with an open
reference implementation. We are not aware of a prior open-source project
structured explicitly around a versioned, contribution-open pattern
catalog of this kind, though we note the field is moving quickly and
invite correction via the project's issue tracker.

\section{A Taxonomy of AI-Associated Vulnerability Patterns}
\label{sec:taxonomy}

We propose the following seven patterns. For each, we state the pattern,
a proposed causal account of why it is disproportionately associated with
AI-assisted code generation specifically (as opposed to being simply
possible in any code), and our detection approach.

\subsection*{AIVP-001: Hardcoded Secret in Example/Placeholder Form}
A plausible-looking example credential (API key, connection string) is
generated to make a code sample runnable, then copy-pasted into real code
without substitution. \emph{Detection}: known-key-format regular
expressions (AWS, Stripe, GitHub token shapes) combined with Shannon
entropy scoring on string literals assigned to secret-like variable names.

\subsection*{AIVP-002: Unparameterized Query Construction}
String interpolation into a query is the more narratively direct way to
express ``build a query using this variable'' than explaining
parameterization. \emph{Detection}: pattern matching on query-execution
calls containing template-literal or f-string interpolation adjacent to a
SQL keyword.

\subsection*{AIVP-003: Shell Command Built by String Interpolation}
Analogous to AIVP-002: \texttt{exec("cmd " + arg)} reads as the intuitive
example; an argument-array-based safe call is correct but a less
frequently generated default. \emph{Detection}: pattern matching on
shell-execution calls with concatenated or interpolated arguments.

\subsection*{AIVP-004: Permissive Default Configuration}
Wildcard CORS origins, disabled TLS verification, and debug flags left
enabled reduce setup friction in a self-contained example, where the
model lacks the real-world context (actual allowed origins, a real
certificate) needed to generate the hardened equivalent.
\emph{Detection}: direct pattern matching on known insecure configuration
flags.

\subsection*{AIVP-005: Weak Cryptographic Primitive in a Security Context}
MD5/SHA1 and non-cryptographic random sources (\texttt{Math.random()})
are frequently the first primitive associated with a generic
``hash this'' or ``generate a random string'' request; the
security-context distinction is not always surfaced absent an explicit
prompt. \emph{Detection}: pattern matching on weak-primitive calls
adjacent to security-relevant identifiers (\texttt{password},
\texttt{token}, \texttt{secret}).

\subsection*{AIVP-006: Inconsistent Authorization Across Near-Identical Routes}
We consider this the most distinctly AI-associated pattern in the set.
When an assistant is asked to add a route ``like the others,'' it
regenerates the handler from a natural-language description rather than
duplicating the existing code block; regeneration is where a
previously-consistent step (an authorization middleware call) can
silently drop, particularly across separate prompts or edits. We propose
this differs structurally from a typical human-introduced authorization
bug, where copy-paste duplication of an existing block tends to
\emph{preserve} the full pattern by construction. \emph{Detection}: for a
given file, identify route-defining calls, determine the middleware
identifier applied to a majority of routes, and flag routes lacking it.

\subsection*{AIVP-007: Verbose Error Response Leaking Internals}
Returning a raw exception object or stack trace directly in an HTTP
response is the fastest way to make error handling visibly ``work'' during
iterative generation; separating server-side logging from a generic
client-facing message is a production concern not always present in a
first-pass generation. \emph{Detection}: pattern matching on error
objects or stack-trace values passed directly into response-sending
calls.

\medskip
\noindent The full catalog, including contribution guidelines for
proposing additional patterns independent of whether a detection rule
yet exists, is maintained openly at the project repository.

\section{Reference Implementation}
\label{sec:tool}

We implement \texttt{ai-vuln-scan}, an open-source (MIT-licensed) command-line
scanner covering AIVP-001 through AIVP-005 and AIVP-007 via regular-expression
and structural pattern matching, and AIVP-006 via a lightweight
route-extraction and majority-middleware heuristic (Section~\ref{sec:taxonomy}).
The tool supports JavaScript/TypeScript and Python, integrates as a CI gate
or pre-commit hook (non-zero exit on any critical-severity finding), and
ships with a labeled example corpus (Section~\ref{sec:evaluation}) usable
as a regression fixture.

We explicitly do not implement full abstract-syntax-tree-based taint
analysis in this version; detection is regex- and structural-heuristic-based,
which we discuss as a limitation in Section~\ref{sec:limitations}.

\section{Evaluation}
\label{sec:evaluation}

\subsection{Corpus}

We constructed a small hand-written test corpus of two files (one
JavaScript/Express, one Python/Flask) deliberately containing instances of
each pattern in Section~\ref{sec:taxonomy}, modeled on realistic
AI-assistant output rather than contrived edge cases. A third file
contains the same routes rewritten to be free of these issues, used to
check for false positives. The full corpus is published alongside the
tool. We are explicit that this is a small, hand-constructed corpus, not
a random or representative sample of real-world AI-generated code; the
evaluation below should be read as a demonstration of feasibility and a
starting comparison point, not a general performance claim.

\subsection{Baselines}

We compare against two established, independently maintained static
analysis tools, run with their default, fully offline configurations:

\begin{itemize}
  \item \textbf{Bandit} \cite{bandit} (v1.9), a widely used Python
  security linter, run against the Python file in our corpus.
  \item \textbf{eslint-plugin-security} \cite{eslintsecurity} (v4.0), a
  widely used JavaScript security rule plugin for ESLint, run against the
  JavaScript file in our corpus with its full rule set enabled.
\end{itemize}

We note that Semgrep \cite{semgrep}, arguably the most prominent tool in
this space, ships without a bundled offline rule set --- its default
usage mode fetches rules from a hosted registry, which was not accessible
in our evaluation environment. We report this as a methodological
limitation rather than omit Semgrep silently; a network-enabled
replication comparing against Semgrep's registry rule sets is noted as
future work.

\subsection{Results}

\begin{table}[h]
\centering
\small
\begin{tabular}{@{}lccc@{}}
\toprule
Pattern & \texttt{ai-vuln-scan} & Bandit & eslint-plugin-security \\
\midrule
AIVP-001 (hardcoded secret) & \checkmark (both files) & \ding{55} & \ding{55} \\
AIVP-002 (unparam. query, direct) & \checkmark (JS) & --- & \ding{55} \\
AIVP-002 (unparam. query, via variable) & \ding{55} (Python, missed) & \checkmark (Python) & --- \\
AIVP-003 (shell interpolation) & \checkmark (both) & \checkmark (Python) & \ding{55} \\
AIVP-004 (permissive defaults) & \checkmark (both) & \checkmark (Python, partial) & \ding{55} \\
AIVP-005 (weak crypto/random) & \checkmark (both) & \checkmark (Python, MD5 only) & \ding{55} \\
AIVP-006 (auth inconsistency) & \checkmark (JS) & --- & \ding{55} (no analogue) \\
AIVP-007 (error leakage) & \checkmark (Python) & \ding{55} & --- \\
False positives (clean file) & 0 & --- & 0 \\
\bottomrule
\end{tabular}
\caption{Detection results by pattern. \checkmark = detected, \ding{55} =
missed, --- = pattern not applicable to that tool's target language in
our corpus. AIVP-002 is split into two rows because our tool's regex-based
approach detects the pattern when the tainted string is passed directly
to a query call but --- a limitation we report candidly --- misses the
case where the same tainted string is first assigned to an intermediate
variable, which Bandit's more sophisticated analysis catches correctly.}
\label{tab:results}
\end{table}

Across our corpus, \texttt{ai-vuln-scan} produced 13 findings (6 critical,
4 high, 3 medium severity) with zero false positives against the clean
comparison file. Bandit produced 7 findings against the Python file,
correctly identifying 4 of 6 patterns we targeted plus one
(unparameterized-query-via-variable) that our tool missed, alongside two
findings (subprocess-import advisory, missing-request-timeout) outside
the scope of our taxonomy entirely. \texttt{eslint-plugin-security},
run with its full default rule set, produced zero findings against our
JavaScript file --- consistent with that plugin's design focus on a
different vulnerability class (unsafe regular expressions, non-literal
filesystem/require calls, timing attacks) rather than the patterns in our
taxonomy.

\subsection{Discussion}

Two results are worth highlighting. First, the AIVP-006
(authorization-inconsistency) pattern has no analogue in either baseline
tool's detection surface --- neither tool reasons about consistency across
multiple route definitions within a file. We view this as modest evidence
that at least one pattern in our taxonomy identifies a genuinely
under-served detection surface, though we caution against generalizing
from a single corpus instance.

Second, and reported candidly rather than omitted: our tool missed a
SQL-injection instance that Bandit correctly caught, because our
detection logic requires the tainted string literal to appear directly
in the query-call argument, and does not trace data flow through an
intermediate variable assignment. This is a direct consequence of using
regex/structural heuristics rather than proper taint analysis, and we
discuss it further in Section~\ref{sec:limitations}. We consider it
important to report this result rather than select only favorable
comparisons, both as good methodological practice and because the honest
gap is itself informative: it demonstrates precisely why a
production-grade version of this tool would need AST-based taint
tracking rather than pattern matching alone.

\section{Limitations}
\label{sec:limitations}

\begin{itemize}
  \item \textbf{Corpus size and construction.} Our evaluation corpus is
  small and hand-constructed by the authors, not sampled from real-world
  AI-generated commits. It should be read as a feasibility demonstration,
  not a claim of general detection performance. A follow-up study using a
  larger corpus of real, labeled AI-assistant commit diffs is the
  clearest next step and is left as future work.
  \item \textbf{Detection approach.} Regex/structural heuristics are
  precise for the patterns they target but will miss variants involving
  intermediate variables, aliasing, or multi-file data flow, as
  demonstrated in Section~\ref{sec:evaluation}. A taint-analysis-based
  successor is planned.
  \item \textbf{Semgrep comparison incomplete.} As noted, we could not
  evaluate against Semgrep's hosted registry rule sets in our
  environment; this comparison is left for a network-enabled replication.
  \item \textbf{Causal claims are proposed, not directly tested.} The
  ``why'' explanations in Section~\ref{sec:taxonomy} are plausible
  accounts consistent with the patterns observed, not the product of a
  controlled study isolating AI-assistant generation behavior against a
  human-authored control group. We regard this as the most important
  open question for follow-up work: a systematic comparison of matched
  human- and AI-authored code samples would allow these causal claims to
  be tested directly rather than proposed.
\end{itemize}

\section{Conclusion}

We present a preliminary, seven-pattern taxonomy of vulnerabilities
disproportionately associated with AI-assistant-generated code, an
open-source reference scanner implementing detection for six of the seven
patterns, and a small empirical comparison against two established
baseline tools. We report our results candidly, including a case where
our tool underperformed an established baseline, and frame the taxonomy
itself --- published as an open, versioned, contribution-welcoming
artifact independent of any particular tool's implementation --- as the
primary contribution of this work. We hope this provides a starting point
for more rigorous, larger-scale study of this problem by the broader
security research community.

\section*{Availability}
The taxonomy, reference implementation, and evaluation corpus described
in this paper are published under the MIT license at
\url{https://github.com/jitendrarout/ai-vuln-scan}.

\begin{thebibliography}{9}

\bibitem{bandit}
PyCQA. \emph{Bandit: A security linter for Python code.}
\url{https://bandit.readthedocs.io}

\bibitem{eslintsecurity}
ESLint Community. \emph{eslint-plugin-security.}
\url{https://github.com/eslint-community/eslint-plugin-security}

\bibitem{semgrep}
Semgrep, Inc. \emph{Semgrep: Lightweight static analysis for many languages.}
\url{https://semgrep.dev}

\end{thebibliography}

\end{document}
